\documentclass[12pt]{article}
\usepackage[T1]{fontenc}
\usepackage[utf8]{inputenc}
\usepackage{lmodern}
\usepackage{textcomp}
\usepackage{enumitem}
\usepackage{geometry}
\geometry{margin=1in}
\begin{document}
\begin{center}
Answer Key - Geography - Undergraduate Level Assessment 20 \
\small (Answers are ordered exactly as in the question paper.)
\end{center}

\section*{Multiple Choice}
\begin{enumerate}[label=\arabic*.]
\item \textbf{Answer:} A
\\ \textit{Options:} A) Colombia, Peru, Uruguay, and Costa Rica.; B) Bhutan, Niger, Cambodia, and Costa Rica.; C) Uruguay, Costa Rica, Timor-Leste, and Pakistan.; D) Bhutan, Cambodia, Uruguay and Costa Rica.
\\ \textit{Note:} Colombia, Peru, Uruguay, and Costa Rica are recognized for their significant investments in education and progressive policies that have led to high secondary school enrollment rates compared to other less developed countries.
\item \textbf{Answer:} C
\\ \textit{Options:} A) 20\%; B) 40\%; C) 60\%; D) 80\%
\\ \textit{Note:} The correct answer of 60\% reflects the significant public support for the incumbent president's decision to seek a third term, indicating a favorable perception of their leadership and policies among the American populace at that time.
\item \textbf{Answer:} C
\\ \textit{Options:} A) 690 million; B) 6.9 billion; C) 69 billion; D) 690 billion
\\ \textit{Note:} The correct answer of 69 billion reflects the staggering scale of global poultry production in 2018, where chickens accounted for the majority of meat consumption worldwide, underscoring their significance in agricultural economics and food supply chains.
\item \textbf{Answer:} B
\\ \textit{Options:} A) The global growth rate was four times as high 50 years ago as it is in 2020.; B) The global growth rate was two times as high 50 years ago as it is in 2020.; C) The global growth rate is two times as high as it is in 2020.; D) The global growth rate is four times as high today as it is in 2020.
\\ \textit{Note:} The correct answer is based on historical demographic data, which shows that the global growth rate in the 1970 s was significantly higher, approximately 2.2\% per year, compared to around 1.0\% in 2020, highlighting a substantial decline in growth over the past five decades.
\item \textbf{Answer:} B
\\ \textit{Options:} A) Japan; B) Canada; C) Russia; D) Iran
\\ \textit{Note:} Canada had one of the highest per capita CO 2 emissions in 2017 due to its reliance on fossil fuels for energy production, extensive industrial activities, and a low population density that results in higher emissions per person compared to many other countries.
\end{enumerate}

\section*{True / False}
\begin{enumerate}[label=\arabic*.]
\setcounter{enumi}{5}
\item \textbf{Answer:} False
\\ \textit{Options:} A) True; B) False
\\ \textit{Note:} The correct answer is: bicameral
\item \textbf{Answer:} True
\\ \textit{Options:} A) True; B) False
\\ \textit{Note:} According to the passage: North Carolina State Fairgrounds
\item \textbf{Answer:} True
\\ \textit{Options:} A) True; B) False
\\ \textit{Note:} According to the passage: two hours
\item \textbf{Answer:} False
\\ \textit{Options:} A) True; B) False
\\ \textit{Note:} The correct answer is: Between 1836 and 1842
\item \textbf{Answer:} True
\\ \textit{Options:} A) True; B) False
\\ \textit{Note:} According to the passage: 10 years
\end{enumerate}

\section*{Long Form Response}
\begin{enumerate}[label=\arabic*.]
\setcounter{enumi}{10}
\item \textbf{Answer:} Venice
\\ \textit{Note:} Expected answer: Venice
\item \textbf{Answer:} coastal region
\\ \textit{Note:} Expected answer: coastal region
\end{enumerate}

\vfill
\noindent\textit{End of Answer Key}
\end{document}